% Source: https://artofproblemsolving.com/community/c3h2151093_the_hardest_math_problem
% Stable identity: aops-topic-2151093-q01
The statement is \textbf{false}: no negative starting value ever reaches $1$.

Indeed, if the current number $m$ is negative, then the next number is either $\dfrac {m}{2}<0$ (if $m$ is even) or $3m+1\le3(-1)+1=-2<0$ (if $m$ is odd). So once a term is negative, every later term is negative, and the sequence never contains $1$.

For example,
\[-1\to-2\to-1\to\cdots,\qquad -5\to-14\to-7\to-20\to-10\to-5\to\cdots\]
are cycles, and $-17$ returns to $-17$ after $18$ steps (passing through $-50,-25,-74,-37,\dots,-272,-136,-68,-34$). So the claim fails for every negative $n$, for example $\boxed{n=-1}$.

\emph{Remark.} For positive starting values the claim is the famous Collatz conjecture, which is still unproved; it has been checked by computer for all starting values up to at least $2^{68}$.
